\documentclass[11pt]{article}

\usepackage[margin=1in]{geometry}
\usepackage{amsmath,amssymb,bm}

% ---------------------------------------------------------------------------
% Notation macros
% ---------------------------------------------------------------------------
\newcommand{\ud}{\,\mathrm{d}}
\newcommand{\dive}{\operatorname{div}}      % divergence w.r.t. spatial coords x
\newcommand{\Dive}{\operatorname{Div}}      % divergence w.r.t. referential coords X
\newcommand{\grad}{\operatorname{grad}}     % gradient  w.r.t. spatial coords x
\newcommand{\Grad}{\operatorname{Grad}}     % gradient  w.r.t. referential coords X
\newcommand{\bx}{\bm{x}}
\newcommand{\bX}{\bm{X}}
\newcommand{\bu}{\bm{u}}
\newcommand{\bv}{\bm{v}}
\newcommand{\bb}{\bm{b}}
\newcommand{\bn}{\bm{n}}
\newcommand{\bN}{\bm{N}}
\newcommand{\bt}{\bm{t}}
\newcommand{\bT}{\bm{T}}
\newcommand{\bw}{\bm{w}}
\newcommand{\bsig}{\bm{\sigma}}
\newcommand{\bP}{\bm{P}}
\newcommand{\bS}{\bm{S}}
\newcommand{\bF}{\bm{F}}
\newcommand{\bC}{\bm{C}}
\newcommand{\bE}{\bm{E}}
\newcommand{\bI}{\bm{I}}
\newcommand{\bphi}{\bm{\varphi}}

\title{Strong and Weak Forms of Solid Mechanics:\\
Spatial, Referential, and Total Lagrangian Formulations}
\author{}
\date{}

\begin{document}
\maketitle

%============================================================================
\section{Kinematics and Notation}
%============================================================================

Let $\Omega_R \subset \mathbb{R}^d$ denote the reference (material)
configuration of a body (referential quantities carry the subscript $R$)
with boundary $\partial\Omega_R = \Gamma_R^D \cup
\Gamma_R^N$, $\Gamma_R^D \cap \Gamma_R^N = \emptyset$, and let $\Omega_t$
denote the current (spatial) configuration at time $t$ with boundary
$\partial\Omega_t = \Gamma_t^D \cup \Gamma_t^N$.  The motion
\begin{equation}
  \bx = \bphi(\bX, t), \qquad \bX \in \Omega_R,
\end{equation}
maps material points $\bX$ to spatial points $\bx$.  The displacement,
deformation gradient, and Jacobian are
\begin{equation}
  \bu(\bX,t) = \bphi(\bX,t) - \bX,
  \qquad
  \bF = \Grad \bphi = \bI + \Grad \bu,
  \qquad
  J = \det \bF > 0,
\end{equation}
where $\Grad(\cdot) = \partial(\cdot)/\partial\bX$ and
$\grad(\cdot) = \partial(\cdot)/\partial\bx$ denote referential and spatial
gradients, respectively.  The right Cauchy--Green tensor and the
Green--Lagrange strain are
\begin{equation}
  \bC = \bF^{\!\top}\bF,
  \qquad
  \bE = \tfrac{1}{2}\left(\bC - \bI\right)
      = \tfrac{1}{2}\left(\Grad\bu + \Grad^{\!\top}\bu
        + \Grad^{\!\top}\bu \, \Grad\bu\right).
\end{equation}
The velocity is $\bv = \dot{\bphi} = \dot{\bu}$ (material time derivative).
Mass conservation relates the referential and spatial densities by
$\rho_R = J\rho$.

The stress measures used below are the Cauchy stress $\bsig$ (force per unit
\emph{current} area), the first Piola--Kirchhoff stress $\bP$ (force per unit
\emph{reference} area), and the second Piola--Kirchhoff stress $\bS$, related
through the Piola transform
\begin{equation}
  \bP = J \bsig \bF^{-\top},
  \qquad
  \bS = \bF^{-1}\bP = J \bF^{-1} \bsig \bF^{-\top},
  \qquad
  \bsig = \tfrac{1}{J}\,\bP\bF^{\!\top} = \tfrac{1}{J}\,\bF\bS\bF^{\!\top}.
  \label{eq:piola}
\end{equation}
Note that $\bS = \bS^{\top}$ and $\bsig = \bsig^{\top}$ (balance of angular
momentum), while $\bP$ is in general unsymmetric, satisfying
$\bP\bF^{\!\top} = \bF\bP^{\top}$.

%============================================================================
\section{Strong Form}
%============================================================================

\subsection{Spatial (Eulerian / updated) form}

Balance of linear momentum on the current configuration $\Omega_t$ reads:
find $\bu$ (equivalently the motion $\bphi$) such that
\begin{subequations}
\label{eq:strong-spatial}
\begin{alignat}{2}
  \rho\,\dot{\bv} &= \dive \bsig + \rho\, \bb
      &\quad& \text{in } \Omega_t, \label{eq:mom-spatial}\\
  \bu &= \bar{\bu} && \text{on } \Gamma_t^D, \\
  \bsig\,\bn &= \bar{\bt} && \text{on } \Gamma_t^N, \\
  \bu(\cdot,0) = \bu_0, \quad \bv(\cdot,0) &= \bv_0 && \text{in } \Omega_R,
\end{alignat}
\end{subequations}
where $\rho$ is the current mass density, $\bb$ the body force per unit mass,
$\bn$ the outward unit normal on $\partial\Omega_t$, $\bar{\bt}$ the
prescribed traction per unit current area, and
$\dot{\bv} = \partial\bv/\partial t + (\grad\bv)\bv$ the material
acceleration.  For quasi-statics the inertial term $\rho\,\dot{\bv}$ is
dropped.  Angular momentum balance requires $\bsig = \bsig^{\top}$.

\subsection{Referential (Lagrangian) form}

Pulling \eqref{eq:strong-spatial} back to the reference configuration
$\Omega_R$ using the Piola transform \eqref{eq:piola} and Nanson's formula
$\bn \ud a = J\bF^{-\top}\bN \ud A$, the momentum balance becomes: find
$\bu : \Omega_R \times [0,T] \to \mathbb{R}^d$ such that
\begin{subequations}
\label{eq:strong-referential}
\begin{alignat}{2}
  \rho_R\, \ddot{\bu} &= \Dive \bP + \rho_R\, \bb
      &\quad& \text{in } \Omega_R, \label{eq:mom-referential}\\
  \bu &= \bar{\bu} && \text{on } \Gamma_R^D, \\
  \bP\,\bN &= \bar{\bT} && \text{on } \Gamma_R^N, \\
  \bu(\cdot,0) = \bu_0, \quad \dot{\bu}(\cdot,0) &= \bv_0 && \text{in } \Omega_R,
\end{alignat}
\end{subequations}
where $\rho_R = J\rho$ is the reference density, $\bN$ the outward unit
normal on $\partial\Omega_R$, and $\bar{\bT} = \bar{\bt}\,(\ud a/\ud A)$ the
prescribed nominal traction (force per unit \emph{reference} area).  All
derivatives are taken with respect to $\bX$ over the fixed domain $\Omega_R$;
this is the natural setting for a total Lagrangian discretization.

The system is closed by a constitutive relation.  For a hyperelastic material
with stored energy density $\Psi(\bF)$ (per unit reference volume),
\begin{equation}
  \bP = \frac{\partial \Psi}{\partial \bF},
  \qquad\text{or equivalently}\qquad
  \bS = 2\,\frac{\partial \Psi}{\partial \bC}
      = \frac{\partial \Psi}{\partial \bE}.
\end{equation}

%============================================================================
\section{Weak Form: Total Lagrangian Formulation}
%============================================================================

Define the trial and test spaces over the \emph{fixed} reference domain,
\begin{equation}
  \mathcal{S}_t = \left\{ \bu \in [H^1(\Omega_R)]^d :
      \bu = \bar{\bu} \text{ on } \Gamma_R^D \right\},
  \qquad
  \mathcal{V} = \left\{ \bw \in [H^1(\Omega_R)]^d :
      \bw = \bm{0} \text{ on } \Gamma_R^D \right\}.
\end{equation}
Taking the inner product of \eqref{eq:mom-referential} with
$\bw \in \mathcal{V}$, integrating over $\Omega_R$, and using integration by
parts together with $\bP\bN = \bar{\bT}$ on $\Gamma_R^N$ and $\bw = \bm{0}$
on $\Gamma_R^D$, yields the total Lagrangian weak form:

\medskip
\noindent\emph{Find $\bu(\cdot,t) \in \mathcal{S}_t$ such that, for all
$\bw \in \mathcal{V}$,}
\begin{equation}
  \boxed{\;
  \int_{\Omega_R} \rho_R\, \ddot{\bu} \cdot \bw \ud V
  + \int_{\Omega_R} \bP(\bF) : \Grad \bw \ud V
  = \int_{\Omega_R} \rho_R\, \bb \cdot \bw \ud V
  + \int_{\Gamma_R^N} \bar{\bT} \cdot \bw \ud A
  \;}
  \label{eq:weak-TL-P}
\end{equation}
with $\bF = \bI + \Grad\bu$.  Every integral is evaluated on the reference
configuration, which never changes during the computation---the defining
feature of the total Lagrangian formulation.

\subsection{Equivalent form in the second Piola--Kirchhoff stress}

Substituting $\bP = \bF\bS$ and noting that
$\bF\bS : \Grad\bw = \bS : \bF^{\!\top}\Grad\bw = \bS : \delta\bE[\bw]$
(the last step using the symmetry of $\bS$), where
\begin{equation}
  \delta\bE[\bw]
  = \tfrac{1}{2}\left( \bF^{\!\top}\Grad\bw + \Grad^{\!\top}\bw\,\bF \right)
\end{equation}
is the variation of the Green--Lagrange strain in the direction $\bw$, the
weak form \eqref{eq:weak-TL-P} may equivalently be written
\begin{equation}
  \int_{\Omega_R} \rho_R\, \ddot{\bu} \cdot \bw \ud V
  + \int_{\Omega_R} \bS(\bE) : \delta\bE[\bw] \ud V
  = \int_{\Omega_R} \rho_R\, \bb \cdot \bw \ud V
  + \int_{\Gamma_R^N} \bar{\bT} \cdot \bw \ud A .
  \label{eq:weak-TL-S}
\end{equation}
This is the form commonly used with $(\bS,\bE)$ constitutive pairs
(e.g.\ St.~Venant--Kirchhoff or models expressed in terms of $\bC$).

\subsection{Remarks}

\begin{itemize}
  \item \textbf{Quasi-statics.} Dropping the inertial term in
    \eqref{eq:weak-TL-P} gives the virtual work statement
    $\int_{\Omega_R} \bP : \Grad\bw \ud V
     = \int_{\Omega_R} \rho_R \bb\cdot\bw \ud V
     + \int_{\Gamma_R^N} \bar\bT\cdot\bw \ud A$,
    i.e.\ internal virtual work equals external virtual work.
  \item \textbf{Nonlinearity and linearization.} The weak form is nonlinear
    in $\bu$ through $\bP(\bF)$ (and geometrically through $\bF$).  A Newton
    scheme requires the consistent tangent
    $\mathbb{A} = \partial\bP/\partial\bF
     = \partial^2\Psi/\partial\bF\,\partial\bF$,
    giving the directional derivative
    $D_{\Delta\bu}\!\int_{\Omega_R} \bP : \Grad\bw \ud V
     = \int_{\Omega_R} \Grad\bw : \mathbb{A} : \Grad\Delta\bu \ud V$;
    in the $(\bS,\bE)$ form this splits into the familiar material and
    geometric stiffness contributions.
  \item \textbf{Follower loads.} A traction prescribed per unit
    \emph{current} area (e.g.\ pressure $\bar\bt = -p\,\bn$) becomes
    configuration-dependent when pulled back:
    $\bar\bT = -p\,J\bF^{-\top}\bN$, and then contributes to the tangent
    stiffness.  The form above assumes $\bar\bT$ is a dead (nominal) load;
    the follower pressure is treated in Section~\ref{sec:follower}.
  \item \textbf{Relation to the spatial weak form.} Pushing
    \eqref{eq:weak-TL-P} forward with
    $\bP:\Grad\bw = J\,\bsig : \grad\bw$ and $\ud v = J \ud V$ recovers the
    updated/spatial weak form
    $\int_{\Omega_t} \rho\, \dot\bv\cdot\bw \ud v
     + \int_{\Omega_t} \bsig : \grad\bw \ud v
     = \int_{\Omega_t} \rho\,\bb\cdot\bw \ud v
     + \int_{\Gamma_t^N} \bar\bt\cdot\bw \ud a$;
    the two statements are identical, differing only in the configuration on
    which integrals are evaluated.
\end{itemize}

%============================================================================
\section{Follower Pressure and Load Scheduling}
\label{sec:follower}
%============================================================================

\subsection{Follower pressure}

A pressure $p$ acting per unit \emph{current} area on $\Gamma_R^p$ exerts the
spatial traction $\bar\bt = -p\,\bn$.  Nanson's relation
$\bn \ud a = J\bF^{-\top}\bN \ud A = \operatorname{cof}\bF\,\bN \ud A$
pulls it back to the reference configuration,
\begin{equation}
  \bar\bT(\bu) = -p\,J\bF^{-\top}\bN = -p\,\operatorname{cof}\bF\,\bN,
\end{equation}
so that the external virtual work of the pressure depends on the
displacement.  Moving it to the internal side, the quasi-static residual of
\eqref{eq:weak-TL-P} becomes
\begin{equation}
  R(\bu;\bw) = \int_{\Omega_R} \bP(\bF) : \Grad\bw \ud V
  + \int_{\Gamma_R^p} p\,(\operatorname{cof}\bF\,\bN)\cdot\bw \ud A
  - \int_{\Omega_R} \rho_R\,\bb\cdot\bw \ud V
  - \int_{\Gamma_R^N} \bar\bT\cdot\bw \ud A .
  \label{eq:residual-follower}
\end{equation}
Its directional derivative in $\Delta\bu$ adds the (in general
non-symmetric) load stiffness
\begin{equation}
  D_{\Delta\bu} \int_{\Gamma_R^p} p\,(\operatorname{cof}\bF\,\bN)\cdot\bw \ud A
  = \int_{\Gamma_R^p} p\,\bigl[
      (\bF^{-\top}:\Grad\Delta\bu)\,J\bF^{-\top}\bN
      - J\bF^{-\top}(\Grad\Delta\bu)^{\top}\bF^{-\top}\bN
    \bigr]\cdot\bw \ud A,
  \label{eq:tangent-follower}
\end{equation}
using $D J = J\,\bF^{-\top}:\Grad\Delta\bu$ and
$D\bF^{-\top} = -\bF^{-\top}(\Grad\Delta\bu)^{\top}\bF^{-\top}$.  Because the
tangent is not symmetric, the linear solver must be a general Krylov method
(GMRES); the conjugate gradient option is refused for inputs with a follower
pressure.

In the code (\texttt{kernels/follower\_pressure.hpp}) the term is a boundary
face integrator of the nonlinear form: on each boundary face of $\Gamma_R^p$
the displacement gradient of the adjacent element gives $\bF$ at the face
quadrature points, $\bN \ud A$ is the area-weighted reference normal of the
face map, and the tangent \eqref{eq:tangent-follower} is obtained by seeding
$\bF$ with dual numbers in the direction of each element displacement
degree of freedom, exactly as the material tangent is.  The same face kernel
serves the displacement block of the mixed $u$--$p$ formulation.  The dead
normal pressure, $\bar\bT = -p\,\bN$ per unit \emph{reference} area, remains
a linear-form term; the two coincide at $\bF = \bI$ and differ at $O(p^2)$
for small loads.  Verification: the tangent against finite differences on
distorted quadrilateral and hexahedral meshes, and the inflation of a
thick-walled incompressible cylinder against the closed form
$P = \mu\,[\ln(B/A) - \ln(b/a) + \tfrac12 (A^2 - a^2)(b^{-2} - a^{-2})]$,
$b^2 = B^2 - A^2 + a^2$ (plane strain, neo-Hookean).

\subsection{Load scheduling}

Quasi-static loading is parametrised by a pseudo-time $t \in [0,1]$.  Every
prescribed displacement, traction, pressure and body force $i$ carries its
own scalar schedule $s_i(t)$ (a ramp over $[t_0, t_1]$, a constant, or a
piecewise-linear table) and enters the residual as $s_i(t)$ times its data;
data may in addition depend on $t$ through an expression $f(\bX, t)$.  The
default ramp $s_i = t$ for every entry recovers the proportional load path,
in which $t < 1$ solves a scaled problem and $t = 1$ the problem of the weak
form.  The stepper advances $t$ through user-defined breakpoints and, when a
Newton solve fails, restores the last converged state and halves the
increment (bisection), returning to the planned breakpoint afterwards.
Because the hyperelastic problem without follower loads is conservative,
the converged state at $t = 1$ does not depend on the path; staged loading
(one load ramped and held, a second one ramped afterwards) is therefore a
test of the implementation rather than a different problem.  Follower
pressures are not conservative, and their solutions can be path dependent
in the presence of instabilities.

\end{document}
